% Generated draft: requires mathematical and source review.
\documentclass[11pt]{article}
\usepackage[margin=25mm]{geometry}
\usepackage{fontspec}
\setmainfont{Latin Modern Roman}
\usepackage{amsmath,amssymb,mathtools,unicode-math}
\setmathfont{Latin Modern Math}
\usepackage{enumitem,needspace,xurl,hyperref}
\setlength{\parindent}{0pt}
\setlength{\parskip}{5pt}
\setlength{\emergencystretch}{3em}
\newcommand{\N}{\mathbb{N}}
\newcommand{\Z}{\mathbb{Z}}
\newcommand{\Q}{\mathbb{Q}}
\newcommand{\R}{\mathbb{R}}
\newcommand{\C}{\mathbb{C}}
\newcommand{\F}{\mathbb{F}}
\newcommand{\A}{\mathbb{A}}
\newcommand{\PP}{\mathbb{P}}
\newcommand{\E}{\mathbb{E}}
\newcommand{\eps}{\varepsilon}
\newcommand{\dd}{\,\mathrm{d}}
\begin{document}
\section*{UnsolvedMath ID: 20003161}
\textbf{Source problem code:} catalogue entry
\subsection{Definitions and mathematical statement}
The original problem asks broadly to compute invariants of cohomogeneity-one manifolds and classify the then-new positive-curvature candidates topologically. Cohomogeneity one means that an appropriate compact Lie group action has one-dimensional orbit space. The source does not list its candidates or their group diagrams. The supplied partial result treats a seven-dimensional family denoted \(Q_k\), whose definition and index range are not included; it does not classify all cohomogeneity-one manifolds or prove positive curvature for this family.

Write \(N=2k+1\), with \(N>0\) in the positive-order convention of the supplied formulas. The aid reports
\[
 \pi_1(Q_k)=0,\qquad\pi_2(Q_k)\cong\Z,
\]
\[
 H^*(Q_k;\Z)\cong
 \Z[x,y]/(Nx^2,x^3,x^2y,y^2),\qquad |x|=2,\quad|y|=5.
\]
Thus the nonzero integral cohomology groups are \(H^0,H^2,H^5,H^7\cong\Z\) and \(H^4\cong\Z/N\Z\), with \(x^2\) its indicated torsion generator and \(xy\) the degree-seven generator. If \(N=1\), the degree-four group is trivial. The aid mentions spin/characteristic data but supplies no formulas for them; no unspecified characteristic-class values are assigned here.

For the torsion linking form in dimension seven, the reported exact denominator is \(N\). In a nonsingular cyclic-form convention this means
\[
 b(x^2,x^2)=u/N\pmod\Z,\qquad\gcd(u,N)=1,
\]
with the numerator \(u\), generator choices, and orientation convention not given. The denominator alone is not the complete linking-form isometry type.

For every integer \(m\geq2\), let \(d=\gcd(m,N)\) and \(\Z_m=\Z/m\Z\). The supplied coefficient package gives
\[
 H^j(Q_k;\Z_m)\cong
 \begin{cases}
 \Z_m,&j=0,2,5,7,\\
 \Z/d\Z,&j=3,4,\\
 0,&\text{otherwise}.
 \end{cases}
\]
For the integral Bockstein from \(0\to\Z\xrightarrow{\ m\ }\Z\to\Z_m\to0\), a generator \(c_m\in H^3(Q_k;\Z_m)\) may be chosen so that
\[
 \beta_m(c_m)=\frac Nd\,x^2\in H^4(Q_k;\Z).
\]
When \(d=1\), both sides are zero under the trivial-group convention. In particular, for every prime \(\ell\) and \(a\geq1\),
\[
 |H^3(Q_k;\Z_{\ell^a})|=\ell^{\min(a,v_\ell(N))}.
\]
This recovers all prime-power divisors of \(N\). Members with different positive values of \(2k+1\) cannot be homotopy equivalent, already by their integral torsion groups; the aid states pairwise distinction for its intended \(Q\)-family. The missing index range must exclude any identification of indices producing the same torsion order before extending that wording literally. No full topological or smooth classification follows solely from these coefficient and denominator data.
\subsection{Short English statement}
The supplied partial result gives integral and finite-coefficient invariants for an undefined-in-record \(Q_k\) candidate family. Its Bockstein recovers every prime-power divisor of \(2k+1\), distinguishing different torsion orders. The candidate definitions, characteristic-class values, and complete linking data remain missing.
\subsection{Sources}
\begin{itemize}
\item Local UnsolvedMath registry, permanent ID \(20003161\), exact source code \texttt{catalogue entry}; canonical \texttt{aim-topology-notes.json}, \texttt{notes[248]}, source tag \texttt{problem}.
\item Problem 36 in M. Kerin's compilation following the September 2007 AIM workshop ``Manifolds with nonnegative sectional curvature'': \url{https://aimath.org/WWN/nnsectcurvature/nnsectcurvature.pdf}. The supplied verification reports checking the official five-page PDF and finding exact wording without reconstruction. This is an aid inspection claim; no PDF was fetched here.
\item The original literature field is empty. No authoritative bibliographic citation, defining group diagram for \(Q_k\), spin/characteristic formulas, or linking numerator is supplied.
\item Research attempt 1: result \texttt{partial\_result}, problem status at run \texttt{partially\_solved}, rigor label \texttt{proved}, novelty confidence \texttt{low}; classification \texttt{PARTIAL-PROGRESS}, registry status \texttt{partially\_solved}. Created and updated 14 August 2026. These are supplied assessments requiring review.
\item Attribution: Ulam AI / UnsolvedMath Contributors, CC BY 4.0. The machine-generated research aid is not peer reviewed and requires independent verification before citation.
\end{itemize}
\end{document}
